\chapter{The Kalman Filter}
\label{ch:ch18}
% Function names for pseudocode are prefixed with "Kalman" to stay unique
% across the book.
\SetKwFunction{KalmanPredict}{KalmanPredict}
\SetKwFunction{KalmanUpdate}{KalmanUpdate}
\SetKwFunction{KalmanTrack}{TrackIntruder}
\SetKwFunction{KalmanInitTwo}{InitFromTwoMeasurements}
\SetKwFunction{KalmanAhead}{PredictAhead}

\begin{objectives}
\begin{itemize}
  \item Write down the linear-Gaussian state-space model of a moving drone, $\state_k=\mat{F}\state_{k-1}+\mat{B}\ctrl_k+\vect{w}_k$ and $\meas_k=\mat{H}\state_k+\vect{v}_k$, and say what each of $\mat{F}$, $\mat{H}$, $\mat{Q}$ and $\mat{R}$ encodes.
  \item Derive the predict and update steps of the Kalman filter from two facts about Gaussians, and explain the Kalman gain as a precision-weighted compromise between prediction and measurement.
  \item Build the constant-velocity model in 2D and 3D with the white-noise-acceleration process noise, derive its $\mat{Q}$, and run the filter by hand for a few steps.
  \item Tune $\mat{Q}$ and $\mat{R}$ with the innovation statistics, gate outliers with the Mahalanobis distance, initialise a track from two measurements and cope with missing measurements.
  \item Propagate the covariance several steps ahead to obtain the growing uncertainty ellipses that the safety layer uses, implement the filter in \python, and verify its accuracy and consistency on a simulation.
\end{itemize}
\end{objectives}

% ---------------------------------------------------------------------
\section{Why this matters for a drone swarm}
\label{sec:ch18-motivation}

The planners of Parts~II and~III know where every drone of the swarm is,
because the swarm reports its own positions. An intruder reports nothing. A
foreign drone that crosses the swarm's airspace is seen only through a sensor
--- a radar, a camera with a depth estimate, a receiver of broadcast position
messages --- that delivers a noisy position a few times per second. Suppose
the sensor reports the intruder's position every $\dt=0.1$~s with an error of
about $1$~m in each coordinate. The local avoidance layer of \cref{ch:ch13}
needs the intruder's \emph{velocity} to build a velocity obstacle, and the
obvious way to obtain it is to difference two consecutive reports. That
divides the difference of two $1$-metre errors by $0.1$~s: the velocity
estimate carries a noise of about $\sqrt2\cdot1/0.1\approx14$~m/s, several
times the speed of the intruder itself. \Cref{fig:ch18-raw} shows what this
looks like. The raw positions jitter by metres from one report to the next,
and the differenced velocity is useless. No avoidance rule can work with that
input: \orca would see a different collision cone every tenth of a second, and
a formation would scatter in response to noise.

The Kalman filter turns this stream of jittery reports into a stable estimate
of the intruder's state --- position \emph{and} velocity --- together with a
covariance matrix that says how far the estimate can be trusted. It does so by
carrying a model of how the intruder moves, predicting where it will be at the
next report, and correcting the prediction by the new measurement, weighting
the two by their uncertainties. In the experiment of \cref{fig:ch18-raw},
produced by the code of this chapter, the filter reduces the position error
from $1.18$~m to $0.52$~m (root mean square) and the velocity error from
$16.8$~m/s to $0.79$~m/s, at the cost of a few small matrix products per
report.

In the hybrid architecture of \cref{ch:ch24} the filter is the core of the
\emph{prediction layer}: it tracks the intruder, hands its state and
covariance to the local safety layer (\orca or DWA, \cref{ch:ch13,ch:ch14})
and to the trajectory predictors of \cref{ch:ch20}, and its covariance,
propagated a few seconds ahead, gives the growing uncertainty ellipses with
which the safety layer inflates the intruder's radius. The filter is the
Week~9 topic of the study plan (\cref{ch:appA}), and the milestone of that
week is precisely that your planner receives a stable estimated state instead
of raw observations.

\begin{figure}[tb]
  \centering
  \inputfigure{ch18/raw-vs-filtered}
  \caption{Why raw measurements are unusable for planning. Left: four seconds
  of the experiment of \cref{sec:ch18-experiment}; the sensor reports the
  intruder's position at $10$~Hz with a noise of $1$~m along $x$ and $0.5$~m
  along $y$ (grey), the true path is dashed and the Kalman estimate is the
  blue line. Right: the velocity $v_x$ obtained by differencing consecutive
  reports (grey) swings by tens of metres per second, while the filter's
  velocity estimate (blue) follows the true velocity (dashed) to within a
  fraction of a metre per second.}
  \label{fig:ch18-raw}
\end{figure}

The chapter proceeds as follows. \Cref{sec:ch18-intuition} gives the idea,
\cref{sec:ch18-problem} the linear-Gaussian model, the two facts about
Gaussians on which the filter rests, and the constant-velocity model of a
drone. \Cref{sec:ch18-algorithm} derives the predict and update steps and
explains the gain; \cref{sec:ch18-example} runs the filter by hand for five
steps, and \cref{sec:ch18-properties} proves what it guarantees.
\Cref{sec:ch18-tuning} covers tuning, gating, missing measurements and track
initialisation, \cref{sec:ch18-prediction} the propagation of uncertainty
into the future, and \cref{sec:ch18-experiment} an experiment on simulated
data. \Cref{sec:ch18-variants,sec:ch18-implementation,sec:ch18-drone} close
with variants, implementation notes and the place of the filter in the drone
system.

% ---------------------------------------------------------------------
\section{The idea in plain words}
\label{sec:ch18-intuition}

You have two sources of information about the intruder, and both are
imperfect. The first is a \emph{motion model}: a drone that was at some
position with some velocity a tenth of a second ago is now, very probably, a
tenth of a second further along that velocity. The model is not exact, because
the intruder may accelerate, so its prediction becomes less certain the longer
you rely on it. The second is the \emph{sensor}, which reports the position
with an error that does not accumulate but is large at every single report.
Neither source alone is good enough; combined, they are.

The Kalman filter\index{Kalman filter} is the combination. It keeps a Gaussian
belief about the state, a mean $\hat{\state}$ and a covariance $\mat{P}$, and
alternates two steps (\cref{fig:ch18-idea}). In the \textbf{predict
step}\index{Kalman filter!predict step} it pushes the belief through the
motion model: the mean moves along the velocity and the covariance grows,
because the model's own noise is added. In the \textbf{update
step}\index{Kalman filter!update step} it compares the new measurement with
the predicted measurement; the difference is the \emph{innovation}, the news
that the measurement brings. The state is corrected by a fraction of the
innovation, and that fraction, the \emph{Kalman gain}, is large when the
prediction is uncertain and the sensor accurate, small in the opposite case.
The covariance shrinks, because two independent pieces of information are
worth more than one. Then the cycle repeats.

The arithmetic behind the fraction is that of a weighted average. If you weigh
an object on two scales with variances $\sigma_1^2$ and $\sigma_2^2$, the best
combination weights each reading by its inverse variance, its
\emph{precision}, and the precision of the result is the sum of the two
precisions. The update step is this rule applied to vectors: the prediction
plays the role of the first reading, the measurement that of the second, and
the matrices $\mat{P}^-$ and $\mat{R}$ play the variances. Everything else in
the chapter is the bookkeeping that keeps the rule exact when the state has
several components, the sensor sees only some of them, and time passes between
readings.

\begin{keyidea}
Keep a Gaussian belief $\Normal(\hat{\state},\mat{P})$ over the state. Predict
by pushing it through the motion model: the mean moves, the covariance grows
by $\mat{Q}$. Update by moving the mean towards the measurement with a gain
that weights prediction and measurement by their precisions: the covariance
shrinks. For a linear model with Gaussian noise this recursion is exact and
optimal --- no estimator, linear or not, has a smaller mean-square error.
\end{keyidea}

\begin{figure}[tb]
  \centering
  \inputfigure{ch18/idea}
  \caption{One cycle of the Kalman filter. Left, in the plane of positions:
  the estimate at step $k-1$ (blue ellipse) is moved by the motion model and
  inflated by the process noise into the prediction (grey, dashed); the
  measurement $\meas_k$ (red cross, with its noise circle) pulls the mean
  towards itself, and the updated estimate (green) is both shifted and
  tighter than the prediction. Right: the same cycle as a loop; the planner
  reads $\hat{\state}_k$ and $\mat{P}_k$ after every update.}
  \label{fig:ch18-idea}
\end{figure}

% ---------------------------------------------------------------------
\section{Problem statement and notation}
\label{sec:ch18-problem}

\subsection{The linear-Gaussian state-space model}

\begin{definition}[Linear-Gaussian state-space model]
\label{def:ch18-model}
A \textbf{linear-Gaussian state-space model}\index{state-space model!linear-Gaussian}
for a state $\state_k\in\R^n$ observed through measurements $\meas_k\in\R^m$ at
steps $k=1,2,\dots$ consists of a \textbf{motion model}\index{motion model}
and a \textbf{measurement model},
\begin{align}
  \state_k &= \mat{F}\state_{k-1}+\mat{B}\ctrl_k+\vect{w}_k, & \vect{w}_k&\sim\Normal(\vect{0},\mat{Q}),\label{eq:ch18-motion}\\
  \meas_k &= \mat{H}\state_k+\vect{v}_k, & \vect{v}_k&\sim\Normal(\vect{0},\mat{R}),\label{eq:ch18-measurement}
\end{align}
with an initial belief $\state_0\sim\Normal(\hat{\state}_0,\mat{P}_0)$.
$\mat{F}\in\R^{n\times n}$ is the \textbf{state transition matrix},
$\mat{B}\in\R^{n\times l}$ the control matrix acting on a known input
$\ctrl_k\in\R^l$ (absent for an intruder, $\mat{B}\ctrl_k=\vect{0}$),
$\mat{H}\in\R^{m\times n}$ the \textbf{measurement matrix}, $\mat{Q}$ the
\textbf{process noise covariance}\index{process noise} and $\mat{R}$ the
\textbf{measurement noise covariance}\index{measurement noise}. The noise
vectors $\vect{w}_1,\vect{w}_2,\dots$, $\vect{v}_1,\vect{v}_2,\dots$ and
$\state_0$ are mutually independent. The matrices may depend on $k$; the index
is dropped when they do not.
\end{definition}

The motion model says that the next state is a linear function of the current
one plus a random disturbance: $\mat{F}$ carries the deterministic part of the
motion (a drone keeps its velocity), $\vect{w}_k$ collects what the model
cannot know (the intruder's accelerations), and $\mat{Q}$ says how large these
unknowns are per step. The measurement model says which linear combinations of
the state the sensor sees ($\mat{H}$ selects them) and how noisy it is
($\mat{R}$). The independence assumption is what makes a recursion possible:
once the current state is known, the past measurements carry no further
information about the future.

\begin{notebox}{Notation: $\mat{F}$ and $\mat{H}$, or $\mat{A}$ and $\mat{C}$?}
This book writes the transition matrix as $\mat{F}$ and the measurement
matrix as $\mat{H}$, the convention of the tracking literature
\cite{barshalom2001estimation}. Control texts, and the model predictive
controller of \cref{ch:ch21}, write the same matrices as $\mat{A}$ and
$\mat{C}$; the double-integrator model of \cref{ch:ch02} uses $\mat{A}$ and
$\mat{B}$; \textcite{thrun2005probabilistic} use $\mat{A}$, $\mat{B}$,
$\mat{C}$, and \textcite{welch1995introduction} use $\mat{A}$ and $\mat{H}$.
Beware also that $\mat{Q}$ and $\mat{R}$ are noise covariances here but cost
weights in \cref{ch:ch21}; the notation table in the front matter lists both
meanings.
\end{notebox}

\subsection{What the filter computes}

\begin{definition}[Filtering, prediction, estimate]
\label{def:ch18-filtering}
Write $\meas_{1:k}=(\meas_1,\dots,\meas_k)$ for the measurements received up
to step $k$. The \textbf{filtering problem}\index{filtering problem} is to
compute, at every step, the posterior density $p(\state_k\mid\meas_{1:k})$.
Under \cref{def:ch18-model} this density is Gaussian, and the filter
represents it by its mean and covariance, the \textbf{estimate}
\begin{equation}
  \hat{\state}_k=\E[\state_k\mid\meas_{1:k}],\qquad
  \mat{P}_k=\Cov(\state_k\mid\meas_{1:k}).
  \label{eq:ch18-estimate}
\end{equation}
The density $p(\state_k\mid\meas_{1:k-1})$, before $\meas_k$ has been used,
is the \textbf{prediction}\index{Kalman filter!prediction}; it is Gaussian
too, with mean $\hat{\state}^-_k$ and covariance $\mat{P}^-_k$. As in the
notation table, the superscript minus marks quantities computed before the
current measurement is used.
\end{definition}

The Bayes filter of \cref{ch:ch02} computes the posterior in two moves: it
propagates the previous posterior through the motion model to obtain the
prediction, and it multiplies the prediction by the likelihood
$p(\meas_k\mid\state_k)$ to obtain the new posterior. For general densities
both moves are integrals without a closed form. For the linear-Gaussian model
they reduce to matrix arithmetic, because of two facts about Gaussians.

\subsection{Two facts about Gaussians}

\begin{lemma}[Affine transformation and conditioning of Gaussians]
\label{thm:ch18-gaussian-facts}
Let $\state\sim\Normal(\vect{\mu},\mat{\Sigma})$ in $\R^n$.
\begin{enumerate}[label=(\roman*)]
  \item If $\vect{y}=\mat{A}\state+\vect{b}+\vect{w}$ with
  $\vect{w}\sim\Normal(\vect{0},\mat{Q})$ independent of $\state$, then
  $\vect{y}\sim\Normal(\mat{A}\vect{\mu}+\vect{b},\;\mat{A}\mat{\Sigma}\mat{A}\T+\mat{Q})$.
  \item If $\vect{a}$ and $\vect{b}$ are jointly Gaussian,
  \[
    \begin{pmatrix}\vect{a}\\ \vect{b}\end{pmatrix}\sim
    \Normal\!\left(\begin{pmatrix}\vect{\mu}_a\\ \vect{\mu}_b\end{pmatrix},
    \begin{pmatrix}\mat{\Sigma}_{aa}&\mat{\Sigma}_{ab}\\ \mat{\Sigma}_{ba}&\mat{\Sigma}_{bb}\end{pmatrix}\right)
  \]
  with $\mat{\Sigma}_{bb}$ invertible, then $\vect{a}$ given $\vect{b}$ is
  Gaussian with
  \begin{equation}
    \E[\vect{a}\mid\vect{b}]=\vect{\mu}_a+\mat{\Sigma}_{ab}\mat{\Sigma}_{bb}\inv(\vect{b}-\vect{\mu}_b),\qquad
    \Cov(\vect{a}\mid\vect{b})=\mat{\Sigma}_{aa}-\mat{\Sigma}_{ab}\mat{\Sigma}_{bb}\inv\mat{\Sigma}_{ba}.
    \label{eq:ch18-conditioning}
  \end{equation}
\end{enumerate}
\end{lemma}
\begin{proof}[Proof sketch]
(i) is the closure of Gaussians under affine maps and under sums of
independent Gaussians, stated in \cref{ch:ch02}: the mean and covariance follow
from the rules for affine maps and from the independence of $\vect{w}$. For
(ii), write the joint density as $\exp(-\tfrac12\phi)$ with $\phi$ the
quadratic form of the block covariance, and expand $\phi$ with the block
inverse of $\mat{\Sigma}$ (the Schur-complement identity of the refresher,
\cref{ch:appB}). Collecting the terms that contain $\vect{a}$ gives a
quadratic form in
$\vect{a}-\vect{\mu}_a-\mat{\Sigma}_{ab}\mat{\Sigma}_{bb}\inv(\vect{b}-\vect{\mu}_b)$
with the precision matrix
$(\mat{\Sigma}_{aa}-\mat{\Sigma}_{ab}\mat{\Sigma}_{bb}\inv\mat{\Sigma}_{ba})\inv$,
plus terms that depend on $\vect{b}$ alone and are absorbed by the normalising
constant. The full computation is in \cref{ch:appB} and in
\textcite[ch.~3]{thrun2005probabilistic}.
\end{proof}

Fact (i) is the predict step in disguise: the motion model
\cref{eq:ch18-motion} is an affine map of $\state_{k-1}$ plus independent
noise. Fact (ii) is the update step: the state and the measurement are jointly
Gaussian, and conditioning on the observed value of $\meas_k$ gives the
posterior. Note that the conditional covariance in \cref{eq:ch18-conditioning}
does not depend on the value of $\vect{b}$: the covariances of the Kalman
filter can be computed before any measurement arrives, a fact that
\cref{sec:ch18-prediction} exploits.

\subsection{Motion models for a drone}
\label{sec:ch18-models}

\begin{definition}[Constant-velocity model]
\label{def:ch18-cv}
The \textbf{constant-velocity (CV) model}\index{constant-velocity model} of a
drone in $d$ dimensions ($d=2$ or $3$) has the state $\state=(\pos,\vel)\in\R^{2d}$,
position and velocity, and, for a sampling interval $\dt$,
\begin{equation}
  \mat{F}=\begin{pmatrix}\mat{I}_d&\dt\,\mat{I}_d\\ \mat{0}&\mat{I}_d\end{pmatrix},\qquad
  \mat{Q}=q\begin{pmatrix}\tfrac{\dt^3}{3}\mat{I}_d&\tfrac{\dt^2}{2}\mat{I}_d\\[3pt] \tfrac{\dt^2}{2}\mat{I}_d&\dt\,\mat{I}_d\end{pmatrix},\qquad
  \mat{H}=\begin{pmatrix}\mat{I}_d&\mat{0}\end{pmatrix}\ \text{ or }\ \mat{H}=\mat{I}_{2d},
  \label{eq:ch18-cv}
\end{equation}
where $q>0$ is the \textbf{white-noise-acceleration intensity}\index{white-noise acceleration}
in $\mathrm{m^2/s^3}$. The first $\mat{H}$ describes a sensor that measures
positions only, the second one that measures positions and velocities;
$\mat{R}$ is then $\sigma_p^2\mat{I}_d$, or
$\diag(\sigma_p^2\mat{I}_d,\sigma_v^2\mat{I}_d)$ for a sensor with independent
errors of standard deviation $\sigma_p$ in position and $\sigma_v$ in velocity.
\end{definition}

$\mat{F}$ is the double-integrator matrix of \cref{ch:ch02} with the
acceleration removed: $\pos_k=\pos_{k-1}+\dt\,\vel_{k-1}$ and
$\vel_k=\vel_{k-1}$. The intruder does accelerate, of course; the model treats
its acceleration as noise, and the shape of $\mat{Q}$ is what makes that
treatment consistent.

\begin{proposition}[Process noise of the constant-velocity model]
\label{thm:ch18-cv-q}
Suppose that between two samples the acceleration of the drone along one axis
is a continuous-time white noise $a(t)$ with $\E[a(t)]=0$ and
$\E[a(t)a(s)]=q\,\delta(t-s)$. Then the disturbance of position and velocity
along that axis over a step of length $\dt$ has covariance
$q\begin{psmallmatrix}\dt^3/3&\dt^2/2\\ \dt^2/2&\dt\end{psmallmatrix}$, and
the disturbances of different axes are independent, which gives the $\mat{Q}$
of \cref{eq:ch18-cv}.
\end{proposition}
\begin{proof}
Integrate the continuous-time model $\dot p=v$, $\dot v=a(t)$ over $[0,\dt]$
along one axis. A unit impulse of acceleration at time $\tau\in[0,\dt]$ changes
the velocity at time $\dt$ by $1$ and the position by the remaining time
$\dt-\tau$, because it acts on the velocity, which then integrates into the
position. Superposing all instants,
\[
  \vect{w}=\begin{pmatrix}w_p\\ w_v\end{pmatrix}
  =\int_0^{\dt}\begin{pmatrix}\dt-\tau\\ 1\end{pmatrix}a(\tau)\,d\tau .
\]
Its covariance is
$\E[\vect{w}\vect{w}\T]=\int_0^{\dt}\!\!\int_0^{\dt}
\begin{psmallmatrix}\dt-\tau\\1\end{psmallmatrix}
\begin{psmallmatrix}\dt-\tau'&1\end{psmallmatrix}\E[a(\tau)a(\tau')]\,d\tau\,d\tau'$,
and the delta function collapses the double integral to a single one. With
$s=\dt-\tau$,
\[
  \E[\vect{w}\vect{w}\T]=q\int_0^{\dt}\begin{pmatrix}s^2&s\\ s&1\end{pmatrix}ds
  =q\begin{pmatrix}\dt^3/3&\dt^2/2\\ \dt^2/2&\dt\end{pmatrix}.
\]
Independence across axes follows from independent accelerations along the
axes. In matrix form, $\vect{w}=\int_0^{\dt}e^{\mat{A}(\dt-\tau)}\mat{G}\,a(\tau)\,d\tau$
with $\mat{A}=\begin{psmallmatrix}0&1\\0&0\end{psmallmatrix}$ and
$\mat{G}=(0,1)\T$, and $e^{\mat{A}s}=\mat{I}+s\mat{A}$ because
$\mat{A}^2=\mat{0}$; the self-test of \code{ch18\_kalman.py} checks the closed
form against this integral numerically.
\end{proof}

Two features of $\mat{Q}$ matter in practice. First, $q$ is a \emph{rate}: it
does not change when the sampling interval changes, because $\mat{Q}$ carries
all the dependence on $\dt$. A sensor that suddenly reports at $20$~Hz instead
of $10$~Hz needs a new $\mat{Q}$ but the same $q$
(\cref{sec:ch18-implementation} returns to this). Second, the off-diagonal
term couples position and velocity: an unknown acceleration moves both, in
the same direction. A diagonal $\mat{Q}$ is a common shortcut and a mild
error; $\mat{Q}=q\,\mat{I}$ regardless of $\dt$ is a common shortcut and a
serious one. As a starting value, if the intruder's accelerations have a
magnitude of about $a_{\max}$ and persist for about $T_a$ seconds, take
$q\approx a_{\max}^2\,T_a$, so that the model's velocity uncertainty after
$T_a$, $\sqrt{q\,T_a}$, matches the velocity change $a_{\max}T_a$ that a real
manoeuvre produces; \cref{sec:ch18-tuning} shows how to refine the value from
data.

\paragraph{The constant-acceleration model.}
When the intruder manoeuvres persistently --- long turns, sustained climbs
--- a model that carries the acceleration as a state reacts faster. The
\textbf{constant-acceleration (CA) model}\index{constant-acceleration model}
has the state $(\pos,\vel,\acc)\in\R^{3d}$,
\[
  \mat{F}=\begin{pmatrix}\mat{I}&\dt\mat{I}&\tfrac{\dt^2}{2}\mat{I}\\ \mat{0}&\mat{I}&\dt\mat{I}\\ \mat{0}&\mat{0}&\mat{I}\end{pmatrix},\qquad
  \mat{Q}=q\begin{pmatrix}\tfrac{\dt^5}{20}\mat{I}&\tfrac{\dt^4}{8}\mat{I}&\tfrac{\dt^3}{6}\mat{I}\\[2pt] \tfrac{\dt^4}{8}\mat{I}&\tfrac{\dt^3}{3}\mat{I}&\tfrac{\dt^2}{2}\mat{I}\\[2pt] \tfrac{\dt^3}{6}\mat{I}&\tfrac{\dt^2}{2}\mat{I}&\dt\mat{I}\end{pmatrix},\qquad
  \mat{H}=\begin{pmatrix}\mat{I}&\mat{0}&\mat{0}\end{pmatrix},
\]
where the noise is now a white-noise \emph{jerk} of intensity $q$ in
$\mathrm{m^2/s^5}$; the derivation is that of \cref{thm:ch18-cv-q} with the
vector $\bigl((\dt-\tau)^2/2,\ \dt-\tau,\ 1\bigr)\T$ in place of
$(\dt-\tau,1)\T$. The price is a larger state with a component, the
acceleration, that the sensor sees only through two integrations, so the
estimate is noisier during straight flight. \textcite{barshalom2001estimation}
treat both models in detail, together with turning models, and \cref{ch:ch19}
compares filters on a turning target. This chapter uses the CV model
throughout; \code{ch18\_kalman.py} builds both.
